\documentclass[11pt,a4paper]{article}
\usepackage{xeCJK}
\usepackage{fontspec}
\usepackage{xcolor}
\usepackage{fancyhdr}
\usepackage{booktabs}
\usepackage{array}
\usepackage{geometry}
\usepackage[protrusion=true]{microtype}
\geometry{margin=1.8cm,top=2cm,headheight=15pt}
\linespread{1.05}

\setmainfont{Baskerville}
\setCJKmainfont{Songti SC}
\newfontfamily{\starfont}{Songti SC}

\definecolor{wrongred}{RGB}{198,40,40}
\definecolor{wrongbg}{RGB}{253,235,233}
\definecolor{rulegray}{RGB}{200,200,200}

\setlength{\emergencystretch}{3em}
\setlength{\fboxsep}{3pt}
\setlength{\parindent}{0pt}

\newcommand{\blank}{\underline{\hspace{2.8cm}}}
\newcommand{\checkmarkbox}{\fbox{\rule{0pt}{0.8ex}\rule{0.8ex}{0pt}}}
\newcommand{\STAR}{\textcolor{wrongred}{\starfont ★}}
% 错答（红色）
\newcommand{\W}[1]{\textcolor{wrongred}{#1}}
% 正确答案（加粗）
\newcommand{\A}[1]{\textbf{#1}}

\pagestyle{fancy}
\fancyhf{}
\fancyhead[L]{\footnotesize 英语单词默写 advantage -- a.m.（错词订正）}
\fancyhead[R]{\footnotesize 赵文瀚}
\fancyfoot[C]{\thepage}
\renewcommand{\headrulewidth}{0.4pt}
\renewcommand{\footrulewidth}{0pt}

\begin{document}

\begin{center}
{\LARGE\bfseries 英语单词默写 advantage -- a.m.}\\[0.25em]
{\large 错词订正与复盘}\\[0.35em]
{\footnotesize 2029 届高一英语词汇手册默写（advantage--a.m.）· 2026-09}
\end{center}

\vspace{-0.2em}
\noindent{\color{rulegray}\rule{\textwidth}{0.8pt}}

\vspace{0.7em}
\noindent\fcolorbox{black!15}{black!4}{%
  \parbox{\dimexpr\linewidth-2\fboxsep-2\fboxrule\relax}{%
    \small 日期：2026-09\qquad 来源：作业默写（词汇手册 advantage--a.m.，共 50 词）\qquad 姓名：赵文瀚（班级 6，学号 42）\\[0.3em]
    原卷得分：\textcolor{wrongred}{\textbf{38}}\,/\enspace 50\qquad
    错词：\textcolor{wrongred}{\textbf{12}}\,/\enspace 50\qquad
    重做日期：\underline{\hspace{2.4cm}}}}

\vspace{0.8em}
\noindent 说明：下表按原卷题号列出被判错的 12 个词。\STAR\ 表示该题不仅写错，还暴露出词形/词义的系统问题（拼写、词性、近义混淆），需优先攻破。

\vspace{0.5em}
\renewcommand{\arraystretch}{1.28}
\noindent{\small\begin{tabular}{@{}p{0.75cm}@{\hspace{0.12cm}}p{3.9cm}@{\hspace{0.12cm}}p{4.4cm}@{\hspace{0.12cm}}p{2.7cm}@{\hspace{0.12cm}}p{1.8cm}@{\hspace{0.12cm}}p{2.35cm}@{\hspace{0.12cm}}c@{}}
\toprule
\textbf{题号} & \textbf{中文释义} & \textbf{我的错答} & \textbf{正确答案} & \textbf{错误类型} & \textbf{二刷} & \textbf{掌握} \\
\midrule
5  & 业余爱好的        & \W{ameteur}            & \A{amateur}            & 拼写错误 & \checkmarkbox 过\,\checkmarkbox 未过 & 1 \\
7  & 私事              & \W{（空白）}            & \A{affair}             & 完全没背到 & \checkmarkbox 过\,\checkmarkbox 未过 & 0 \\
9  & 感情              & \W{（空白）}            & \A{affection}          & 完全没背到 & \checkmarkbox 过\,\checkmarkbox 未过 & 0 \\
17 & 积极进取的        & \W{agressive（单 g）}   & \A{aggressive}         & 拼写错误 & \checkmarkbox 过\,\checkmarkbox 未过 & 1 \\
18 & 预料 \STAR        & \W{alter（串了 36 题）} & \A{anticipate}         & 串行 & \checkmarkbox 过\,\checkmarkbox 未过 & 1 \\
19 & 令人愉快的 \STAR  & \W{amusing}            & \A{agreeable}          & 词义混淆 & \checkmarkbox 过\,\checkmarkbox 未过 & 1 \\
27 & 使惊恐 \STAR      & \W{afraid（写成形容词）} & \A{alarm}              & 词形错误 & \checkmarkbox 过\,\checkmarkbox 未过 & 1 \\
35 & 沿着 \STAR        & \W{alongside}          & \A{along}              & 形近混淆 & \checkmarkbox 过\,\checkmarkbox 未过 & 1 \\
36 & （使）变样        & \W{（空白）}            & \A{alter}              & 完全没背到 & \checkmarkbox 过\,\checkmarkbox 未过 & 0 \\
41 & 契约              & \W{（空白）}            & \A{agreement}          & 完全没背到 & \checkmarkbox 过\,\checkmarkbox 未过 & 0 \\
45 & 合计 \STAR        & \W{all}                & \A{add up}             & 词义混淆 & \checkmarkbox 过\,\checkmarkbox 未过 & 1 \\
49 & 宣布者            & \W{（空白）}            & \A{announcer}          & 完全没背到 & \checkmarkbox 过\,\checkmarkbox 未过 & 0 \\
\bottomrule
\end{tabular}}

\vspace{0.4em}
\noindent{\footnotesize 注：错答栏由原卷扫描件比对所得，其中第 4、32 题原卷先写错后自行划改，最终判定为对，未计入错词；第 18 题所写的 \emph{alter} 实为第 36 题答案，属串行；第 41、49 题原卷为空。重做前请对照 {\xeCJKsetup{CJKecglue={}}\texttt{vocab\_advantage\_am\_01\_原卷\_1.jpg}}、{\xeCJKsetup{CJKecglue={}}\texttt{\_2.jpg}} 核对。}

\newpage
\begin{center}
{\Large\bfseries 错因归类与复盘}
\end{center}
\vspace{-0.4em}
\noindent{\color{rulegray}\rule{\textwidth}{0.8pt}}

\vspace{0.5em}
\noindent 说明：按错误性质归类 12 个错词，供汇总对账；逐词二刷结果见第 1 页错词表「二刷」列。

\vspace{0.5em}
\renewcommand{\arraystretch}{1.35}
\begin{center}
\begin{tabular}{@{}p{4.6cm} p{10.3cm} c@{}}
\toprule
\textbf{错误类型} & \textbf{题号} & \textbf{词数} \\
\midrule
完全没背到（空白） & {\small 7, 9, 36, 41, 49} & 5 \\
拼写错误 & {\small 5, 17} & 2 \\
词形错误（词性／单复数） & {\small 27} & 1 \\
词义混淆（写成他词） & {\small 19, 45} & 2 \\
形近混淆（写成形近词） & {\small 35} & 1 \\
串行（写了别题答案） & {\small 18} & 1 \\
\midrule
\textbf{合计} & \textbf{12 题} & \textbf{12} \\
\bottomrule
\end{tabular}
\end{center}

\vspace{0.5em}
\noindent\textbf{关联知识点标签}
\vspace{0.25em}

\noindent\begin{tabular}{@{}p{8.2cm}@{\hspace{0.4cm}}p{8.2cm}@{}}
\begin{minipage}[t]{\linewidth}
\small
\STAR\ \textbf{双写字母}：aggressive（双 g）、amateur（易写成 ameteur）\\[0.2em]
\STAR\ \textbf{近义辨析}：along（沿着）/ alongside（在……旁边）\\[0.2em]
\STAR\ \textbf{形容词 vs 名词}：afraid（害怕的）/ alarm（使惊恐；警报）
\end{minipage} &
\begin{minipage}[t]{\linewidth}
\small
\STAR\ \textbf{词义混淆}：alter（改变）/ anticipate（预料）\\[0.2em]
\STAR\ \textbf{程度与搭配}：add up（合计）/ all（全部）\\[0.2em]
\STAR\ \textbf{范围}：A 字母中段（ad- / af- / ag- / al- / am- / an- 词族）
\end{minipage} \\
\end{tabular}

\vspace{0.65em}
\noindent\fcolorbox{black!15}{black!4}{%
  \parbox{\dimexpr\linewidth-2\fboxsep-2\fboxrule\relax}{%
    \textbf{复盘记录}\\[0.35em]
    最薄弱的板块：\underline{\hspace{9.0cm}}\\[0.55em]
    主要错误原因：\checkmarkbox 完全没背到\quad \checkmarkbox 词形记错\quad \checkmarkbox 拼写不牢\quad \checkmarkbox 词义混淆\quad \checkmarkbox 词组漏词\\[0.55em]
    本次复习的改进措施：\underline{\hspace{10.1cm}}\\[0.55em]
    \hspace{3.2em}\underline{\hspace{10.1cm}}\\[0.55em]
    下次复习日期：\underline{\hspace{2.2cm}}\qquad
    当前掌握度：\checkmarkbox 仍薄弱\quad \checkmarkbox 基本掌握\quad \checkmarkbox 已掌握
  }}

\end{document}
